\documentclass[11pt]{article}
\usepackage[margin=0.8in]{geometry}
\usepackage{amsmath, amssymb}
\usepackage{booktabs}
\usepackage{array}
\usepackage{enumitem}
\usepackage{fancyhdr}
\usepackage{parskip}
\usepackage{bm}

% Header and Footer
\pagestyle{fancy}
\setlength{\headheight}{13.6pt}
\fancyhf{}
\fancyhead[C]{\textbf{Chapter 4: Commonly Used Distributions Summary Notes}}
\fancyfoot[C]{CEE 305 - Spring 2026}
\renewcommand{\headrulewidth}{1pt}
\renewcommand{\footrulewidth}{0.5pt}

\begin{document}

% ==========================================
\section*{Introduction to Probability Distributions}
% ==========================================
This chapter covers the most commonly used probability distributions in statistics and engineering applications. Distributions are divided into two main categories: \textbf{discrete} and \textbf{continuous} random variables.

\subsection*{Discrete Random Variables}
Random variables that take on only a finite or countable number of values.

\textbf{Three Key Discrete Distributions:}
\begin{enumerate}[nosep]
    \item Binomial Distribution
    \item Poisson Distribution
    \item Hyper-geometric Distribution (not covered in detail here)
\end{enumerate}
All of these build upon the foundation of the \textbf{Bernoulli random variable}.

% ==========================================
\section*{Section 4.1: The Bernoulli Distribution}
% ==========================================
\textbf{Definition:} A Bernoulli trial is an experiment with exactly two possible outcomes:
\begin{itemize}[nosep]
    \item \textbf{Success} (labeled as 1) with probability \(p\)
    \item \textbf{Failure} (labeled as 0) with probability \(1-p\)
\end{itemize}

\textbf{Notation:} \(X \sim \text{Bernoulli}(p)\)

\textbf{Probability Mass Function (PMF):}
\[
p(0) = P(X=0) = 1-p,\qquad p(1)=P(X=1)=p,\qquad p(x)=0 \text{ otherwise}
\]

\textbf{Mean and Variance:}
\[
\mu_X = p,\qquad \sigma_X^2 = p(1-p)
\]

\textbf{Examples:} Coin toss (heads = success), selecting a defective component from a population.

% ==========================================
\section*{Section 4.2: The Binomial Distribution}
% ==========================================
\textbf{Definition:} A binomial random variable counts the number of successes in \(n\) independent Bernoulli trials.

\textbf{Conditions:}
\begin{enumerate}[nosep]
    \item Fixed number of trials (\(n\) identical trials)
    \item Two outcomes per trial (success/failure)
    \item Constant probability of success \(p\)
    \item Trials are independent
    \item \(X\) = number of successes in \(n\) trials
\end{enumerate}

\textbf{Notation:} \(X \sim \text{Bin}(n,p)\)

\textbf{Probability Mass Function:}
\[
P(X=x) = \binom{n}{x} p^x (1-p)^{n-x},\quad x=0,1,\dots,n
\]
where \(\binom{n}{x} = \frac{n!}{x!(n-x)!}\).

\textbf{Mean and Variance:}
\[
\mu_X = np,\qquad \sigma_X^2 = np(1-p) = npq\;\;(q=1-p)
\]

\textbf{Important Property:} A binomial random variable is the sum of \(n\) independent Bernoulli(\(p\)) variables: \(X = Y_1 + \cdots + Y_n\).

\textbf{Sampling Rule of Thumb:} The binomial distribution can be used when sampling from a finite population \textbf{if the sample size is no more than 5\% of the population}.

\subsection*{Sample Proportion}
If \(X \sim \text{Bin}(n,p)\), then the sample proportion is \(\hat{p} = X/n\).
\begin{itemize}[nosep]
    \item \textbf{Unbiased:} \(\mu_{\hat{p}} = p\)
    \item \textbf{Standard error:} \(\sigma_{\hat{p}} = \sqrt{p(1-p)/n}\)
    \item \textbf{Estimated standard error:} \(\hat{\sigma}_{\hat{p}} = \sqrt{\hat{p}(1-\hat{p})/n}\)
\end{itemize}

% ==========================================
\section*{Section 4.3: The Poisson Distribution}
% ==========================================
\textbf{Definition:} Models the number of occurrences of an event in a fixed unit of time or space.

\textbf{Applications:}
\begin{itemize}[nosep]
    \item Number of phone calls per hour
    \item Number of machine breakdowns per day
    \item Number of traffic accidents per week
    \item Number of particles in a volume of liquid
\end{itemize}

\textbf{Key Concept:} The Poisson distribution approximates the binomial when \(n\) is large, \(p\) is small, and \(\lambda = np\) is moderate.

\textbf{Notation:} \(X \sim \text{Poisson}(\lambda)\), \(\lambda > 0\) (rate parameter).

\textbf{Probability Mass Function:}
\[
P(X=x) = \frac{e^{-\lambda}\lambda^x}{x!},\quad x=0,1,2,\dots
\]

\textbf{Mean and Variance:}
\[
\mu_X = \lambda,\qquad \sigma_X^2 = \lambda
\]
(Unique property: mean equals variance.)

\subsection*{Estimating the Rate Parameter}
If events are observed for \(t\) units, and \(X \sim \text{Poisson}(\lambda t)\), then:
\[
\hat{\lambda} = \frac{X}{t},\qquad \mu_{\hat{\lambda}} = \lambda,\qquad \sigma_{\hat{\lambda}} = \sqrt{\frac{\lambda}{t}},\qquad \hat{\sigma}_{\hat{\lambda}} = \sqrt{\frac{\hat{\lambda}}{t}}
\]

% ==========================================
\section*{Section 4.5: The Normal Distribution}
% ==========================================
\textbf{Definition:} The normal (Gaussian) distribution is symmetric and bell-shaped. It is the most important continuous distribution in statistics.

\textbf{Notation:} \(X \sim N(\mu,\sigma^2)\), where \(\mu\) is the mean (any real) and \(\sigma^2>0\) is the variance.

\textbf{Probability Density Function (PDF):}
\[
f(x) = \frac{1}{\sigma\sqrt{2\pi}} e^{-(x-\mu)^2/(2\sigma^2)},\quad -\infty < x < \infty
\]

\textbf{Mean and Variance:}
\[
\mu_X = \mu,\qquad \sigma_X^2 = \sigma^2
\]

\textbf{The 68-95-99.7\% Rule (Empirical Rule):}
\begin{itemize}[nosep]
    \item 68\% of data within \(\mu \pm \sigma\)
    \item 95\% of data within \(\mu \pm 2\sigma\)
    \item 99.7\% of data within \(\mu \pm 3\sigma\)
\end{itemize}

\subsection*{The Standard Normal Distribution}
Convert to \(z\)-scores: \(z = \frac{x-\mu}{\sigma}\). Then \(Z \sim N(0,1)\).

\textbf{PDF of Standard Normal:}
\[
f(z) = \frac{1}{\sqrt{2\pi}} e^{-z^2/2},\quad -\infty < z < \infty
\]

\textbf{Properties:} Symmetric about 0, mean 0, variance 1, total area = 1.

\textbf{Finding Probabilities Using the z-Table (Table A.2):}
\begin{enumerate}[nosep]
    \item Convert \(x\) to \(z\)
    \item Find area to left of \(z\) in the table
    \item Area to right = \(1 - \text{area left}\)
    \item Area between \(z_1\) and \(z_2\) = \(\text{area left of }z_2 - \text{area left of }z_1\)
\end{enumerate}

\textbf{Working Backwards (Percentiles):} To find the \(z\)-score for the \(p\)th percentile, locate area \(p\) in the table body and read the corresponding \(z\).

% ==========================================
\section*{Linear Combinations of Normal Random Variables}
% ==========================================
\textbf{Single Linear Transformation:}
If \(X \sim N(\mu,\sigma^2)\) and \(a \neq 0, b\) are constants:
\[
aX + b \sim N(a\mu + b,\; a^2\sigma^2)
\]

\textbf{Linear Combinations:}
If \(X_1,\dots,X_n\) are independent normal, then for constants \(c_i\):
\[
\sum c_i X_i \sim N\left(\sum c_i\mu_i,\; \sum c_i^2\sigma_i^2\right)
\]

\textbf{Important Special Cases:}
\begin{itemize}[nosep]
    \item Sample mean: \(\bar{X} \sim N(\mu,\; \sigma^2/n)\)
    \item Sum of two independent normals: \(X+Y \sim N(\mu_X+\mu_Y,\; \sigma_X^2+\sigma_Y^2)\)
    \item Difference of two independent normals: \(X-Y \sim N(\mu_X-\mu_Y,\; \sigma_X^2+\sigma_Y^2)\)
\end{itemize}

% ==========================================
\section*{Additional Key Concepts}
% ==========================================
\subsection*{Probability Plots}
Used to assess distributional fit.
\begin{enumerate}[nosep]
    \item Sort data in ascending order
    \item Assign probabilities \((i-0.5)/n\) to each value
    \item Plot observed values against theoretical quantiles
    \item If approximately linear, the distribution is a good fit
\end{enumerate}

\subsection*{Central Limit Theorem (CLT)}
Let \(X_1,\dots,X_n\) be a random sample from any population with mean \(\mu\) and variance \(\sigma^2\). For large \(n\):
\[
\bar{X} \approx N\left(\mu,\frac{\sigma^2}{n}\right),\qquad S_n = \sum X_i \approx N(n\mu, n\sigma^2)
\]
\textbf{Rule of Thumb:} CLT is good for \(n>30\) for most populations.

\subsection*{Normal Approximations}
\begin{itemize}[nosep]
    \item \textbf{To Binomial:} If \(X\sim\text{Bin}(n,p)\) with \(np>10\) and \(n(1-p)>10\):
    \[
    X \approx N(np, np(1-p)),\qquad \hat{p} \approx N\left(p,\frac{p(1-p)}{n}\right)
    \]
    \item \textbf{To Poisson:} If \(X\sim\text{Poisson}(\lambda)\) with \(\lambda>10\):
    \[
    X \approx N(\lambda, \lambda)
    \]
\end{itemize}

\subsection*{Simulation}
Process of generating random numbers to model real phenomena:
\begin{enumerate}[nosep]
    \item Generate random samples from specified distributions
    \item Count occurrences of events of interest
    \item Estimate probabilities by relative frequencies
\end{enumerate}

% ==========================================
\section*{Summary Table of Distributions}
% ==========================================
{\small
\begin{center}
\begin{tabular}{l c c c c c}
\toprule
\textbf{Distribution} & \textbf{Type} & \textbf{Parameters} & \textbf{PMF/PDF} & \textbf{Mean} & \textbf{Variance} \\
\midrule
Bernoulli & Discrete & \(p\) & \(P(X=x)=p^x(1-p)^{1-x}\) & \(p\) & \(p(1-p)\) \\
Binomial & Discrete & \(n,p\) & \(P(X=x)=\binom{n}{x}p^x(1-p)^{n-x}\) & \(np\) & \(np(1-p)\) \\
Poisson & Discrete & \(\lambda\) & \(P(X=x)=\frac{e^{-\lambda}\lambda^x}{x!}\) & \(\lambda\) & \(\lambda\) \\
Normal & Continuous & \(\mu,\sigma^2\) & \(f(x)=\frac{1}{\sigma\sqrt{2\pi}}e^{-(x-\mu)^2/(2\sigma^2)}\) & \(\mu\) & \(\sigma^2\) \\
Standard Normal & Continuous & \(0,1\) & \(f(z)=\frac{1}{\sqrt{2\pi}}e^{-z^2/2}\) & \(0\) & \(1\) \\
\bottomrule
\end{tabular}
\end{center}
}

% ==========================================
\section*{Key Terminology}
% ==========================================
\begin{itemize}[nosep]
    \item \textbf{Bernoulli Trial:} An experiment with two possible outcomes (success/failure).
    \item \textbf{Binomial Distribution:} Distribution of the number of successes in \(n\) independent Bernoulli trials.
    \item \textbf{Poisson Distribution:} Distribution of the number of events in a fixed interval of time or space.
    \item \textbf{Normal Distribution:} Continuous, symmetric, bell-shaped distribution.
    \item \textbf{Standard Normal Distribution:} Normal distribution with mean 0 and variance 1.
    \item \textbf{z-score:} Number of standard deviations an observation is from the mean.
    \item \textbf{Central Limit Theorem (CLT):} The sampling distribution of the sample mean approaches normality as sample size increases.
    \item \textbf{Probability Plot:} Graphical tool to assess whether a sample comes from a specific distribution.
    \item \textbf{Simulation:} Using random number generation to model and analyze real-world processes.
\end{itemize}

% ==========================================
\section*{Important Notes}
% ==========================================
\begin{itemize}[nosep]
    \item The Poisson approximation to the binomial works well when \(n\) is large, \(p\) is small, and \(np\) is moderate (e.g., \(n \ge 20\), \(p \le 0.05\)).
    \item The normal approximation to the binomial is good when both \(np\) and \(n(1-p)\) exceed 10.
    \item The Central Limit Theorem applies to any population, regardless of its distribution, provided the sample size is sufficiently large.
    \item For normal probability plots, strong curvature indicates departure from normality; outliers appear as points far from the line.
    \item The mean and variance of the sample proportion are derived directly from the binomial distribution.
\end{itemize}

\end{document}